\section{Actual pure-cube selectors and a proved squarefree-value input}
\label{cu19-section}

The following construction imposes the marked congruences on the root
before cubing. A published unconditional theorem supplies the complete
squarefree-value asymptotic after its hypotheses have been checked on
the resulting integer polynomial. Its constants depend on the fixed
packet; no assertion for changing exponents is made.

Work in \(\mathcal O=\mathbb Z[\zeta]\), where
\(\zeta^2-\zeta+1=0\). For a positive integer \(k\), put \(a=3k\) and
\[
 (a+\zeta)^3=A+C\zeta,\qquad c=A+C,\qquad T_3(k)=ACc.
\]

\begin{proposition}[The actual cube and every irreducible factor]
\label{cu19-actual}
The three arms are positive and pairwise coprime, with
\[
 A=a^3-3a-1,\qquad C=3a(a+1),\qquad c=a^3+3a^2-1.
\]
In particular,
\begin{equation}
 T_3(k)=3a(a+1)f(a)g(a),\qquad
 f(a)=a^3-3a-1,\quad g(a)=a^3+3a^2-1.
 \label{cu19-boundary}
\end{equation}
Both cubics are irreducible over \(\mathbb Q\), each of discriminant
\(81\), and \(a(a+1)f(a)g(a)\) is squarefree over \(\mathbb Q\).
Modulo every prime \(p>3\) this product has no repeated root over an
algebraic closure and has at most eight roots in \(\mathbb F_p\).
Every prime dividing the full boundary is a unit for the actual input
norm \(a^2+a+1\), which is \(1\pmod3\).
\end{proposition}
\begin{proof}
Direct multiplication gives the formulas and positivity for \(a\ge3\).
The values of \(f\) at zero and minus one are respectively \(-1,1\);
also \(f(a)\equiv-1\pmod3\) for \(a=3k\).
Thus \(A\) is coprime to \(3a(a+1)=C\). The other two gcds follow
by taking their sum \(c\). Norm multiplicativity gives
\[
 (a^2+a+1)^3=A^2+AC+C^2.
\]
At a prime dividing one of the three coprime arms, the right side
is a nonzero square modulo that prime. This proves the norm-unit
assertion for every boundary prime.

The rational-root test at \(1,-1\) proves irreducibility of both
monic cubics. The cubic discriminant formula gives \(81\) for each.
Moreover \(g-f=3a(a+1)\), while \(f(0)=-1\) and \(f(-1)=1\).
Outside characteristic three, these identities exclude common
roots between the cubics or between a cubic and a linear factor.
The discriminants exclude repeated roots inside either cubic.
They also prove the corresponding assertions over \(\mathbb Q\).
The substitution \(a=3k\) preserves these conclusions outside three.
\end{proof}

\begin{theorem}[Exact packets and the actual quotient polynomial]
\label{cu19-packet}
Fix a finite nonempty set \(\mathcal S\) of primes \(q>7\). For each
\(q\), fix an actual positive owner \(k(q)\) with
\(e_q=v_q(T_3(k(q)))\ge1\). Choose a positive integer \(k_0\) with
\[
 k_0\equiv k(q)\pmod{q^{e_q+1}}\qquad(q\in\mathcal S).
\]
Define
\[
 \mathcal E=\mathcal S\cup\{2,3,5,7\},\qquad
 F_p=v_p(T_3(k_0)),\qquad
 K=\prod_{p\in\mathcal E}p^{F_p},\qquad
 H=\prod_{p\in\mathcal E}p^{F_p+1}.
\]
For \(j\ge1\) integral put \(k_j=k_0+Hj\). The outputs
\((3k_j+\zeta)^3\) remain positive primitive unramified pure cubes,
and their exact depths at all \(p\in\mathcal E\) are \(F_p\).
In particular the marked depths remain \(e_q\).

The literal polynomial
\begin{equation}
 G(J)=T_3(k_0+HJ)/K
 \label{cu19-quotient}
\end{equation}
belongs to \(\mathbb Z[J]\). It is nonconstant and squarefree over
\(\mathbb Q\), each irreducible factor has degree at most three,
and it has no fixed prime divisor.
\end{theorem}
\begin{proof}
CRT applies to the pairwise coprime prime powers. An integer
polynomial preserves its value modulo these powers, so the extra
digit gives \(F_q=e_q\). Since \(H\) contains \(p^{F_p+1}\) for
every \(p\in\mathcal E\), all these exact valuations remain fixed
along the progression. Proposition~\ref{cu19-actual} proves the
actual power, gcd and norm assertions.

To prove the stronger assertion \(G\in\mathbb Z[J]\), its constant
coefficient is divisible by \(K\), and every nonconstant coefficient
of \(T_3(k_0+HJ)\) is divisible by \(H\). This follows by expanding
each integer monomial. Since \(K\mid H\), division by \(K\) is
coefficientwise integral.

For \(p\in\mathcal E\), the constant \(T_3(k_0)/K\) is a \(p\)-unit,
whereas all nonconstant coefficients after division are multiples
of \(H/K\), which is divisible by \(p\). Thus \(G\) is a fixed
nonzero constant modulo \(p\). For \(p\notin\mathcal E\), the
quantities \(3,H,K\) are units and \(p\ge11\). The affine substitution
is invertible; Proposition~\ref{cu19-actual} gives a nonzero
squarefree polynomial of degree eight with at most eight roots
modulo \(p\). It cannot vanish at every residue class. This excludes
a fixed divisor at every prime. A nonzero affine substitution and
a nonzero scalar preserve factor degrees and squarefreeness over
\(\mathbb Q\), proving the remaining claims.
\end{proof}

\begin{theorem}[Positive density with the complete signed budget]
\label{cu19-density}
Keep all data of Theorem~\ref{cu19-packet} fixed. There is a set
\(\mathcal G\) of positive integers of natural density
\[
 d(\mathcal G)=\prod_p\left(1-\frac{\rho_G(p^2)}{p^2}\right)
 \ge\frac15,\qquad
 \rho_G(m)=\#\{j\bmod m:G(j)\equiv0\pmod m\},
\]
such that each \(j\in\mathcal G\) has
\[
 T_3(k_j)=KR_j,\qquad R_j\text{ squarefree},\qquad
 \gcd(K,R_j)=1,
\]
and the prime divisors of \(R_j\) all lie outside \(\mathcal E\).
For \(J(T)=\sum_{p\mid T}(v_p(T)-3)\log p\), the exact identity is
\begin{equation}
 J(T_3(k_j))=-2\log T_3(k_j)+3\log\frac K{\operatorname{rad}K}.
 \label{cu19-signed}
\end{equation}
Along these good indices tending to infinity,
\[
 \frac{J(T_3(k_j))}{\log c_{k_j}}\longrightarrow-\frac{16}{3},
 \qquad
 \frac{\log\operatorname{rad}(T_3(k_j))}{\log c_{k_j}}
 \longrightarrow\frac83.
\]
Every good output satisfies the explicit bound
\begin{equation}
 c_{k_j}\le
 2\left(\frac K{\operatorname{rad}K}\right)^{3/8}
 \operatorname{rad}(T_3(k_j))^{3/8}.
 \label{cu19-radical}
\end{equation}
\end{theorem}
\begin{proof}
The precise external input is Booker--Browning
\cite[Theorem~1.2, printed p.~4]{BookerBrowning2016Squarefree}.
For a nonconstant squarefree integer polynomial with no fixed
prime divisor and irreducible factors of degree at most three,
their theorem gives
\begin{equation}
 \sum_{1\le j\le X}\mu^2(G(j))
 =X\prod_p\left(1-\frac{\rho_G(p^2)}{p^2}\right)
   +O_G(X/\log X).
 \label{cu19-bb}
\end{equation}
The no-fixed-prime assumption is imposed in the authors' opening
setup on printed p.~1; their equation~(2) defines \(\rho\).
All these hypotheses were individually checked in
Theorem~\ref{cu19-packet}. The theorem is unconditional; its
large-square input is proved using the low-degree results described
in their Section~4. We do not substitute an elementary degree-eight
union estimate for that input, and no ABC conjecture is assumed.

For \(p\in\mathcal E\), the polynomial is always a \(p\)-unit, so
\(\rho_G(p^2)=0\). Outside \(\mathcal E\), its at most eight simple
roots lift uniquely to \(p^2\), so \(\rho_G(p^2)\le8\). Hence
\[
 \sum_{p\notin\mathcal E}\frac{\rho_G(p^2)}{p^2}
 \le8\sum_{m\ge11}m^{-2}
 <8\int_{10}^{\infty}x^{-2}\,dx=\frac45.
\]
The product inequality \(\prod(1-u_i)\ge1-\sum u_i\), followed
by a convergent product limit, proves the lower bound \(1/5\).
Take \(\mathcal G\) to be the indices for which \(G(j)\) is squarefree.
All original factors are positive for \(j\ge1\); exact freezing
excludes \(\mathcal E\) from \(G(j)\). Thus \(R_j=G(j)\) has all
the asserted properties.

Now
\[
 \operatorname{rad}(T_3(k_j))
 =T_3(k_j)/(K/\operatorname{rad}K).
\]
This proves \eqref{cu19-signed}, retaining every frozen depth and
every unmarked negative contribution. The polynomial degrees give
\(\log T_3=8\log j+O(1)\) and \(\log c=3\log j+O(1)\), which prove
the limits. For \(a\ge3\),
\[
 f(a)\ge a^3/2,\qquad g(a)\ge a^3,\qquad 3a(a+1)\ge3a^2.
\]
The first inequality follows from \(a^3-6a-2\ge0\), increasing on
\([3,\infty)\). Therefore \(T_3\ge a^8\), while
\(c=a^3+3a^2-1\le2a^3\). Eliminating \(a\) and using the radical
identity proves \eqref{cu19-radical}.
\end{proof}

\paragraph{Actual distinct owners and arbitrary fixed positive depths.}
At \(k=1,2\), respectively, the literal triples are
\[
 (A,C,c)=(17,36,53),\qquad (197,126,323).
\]
The integers \(53\) and \(197\) are prime by trial division up to
their square roots. The first triple has no factor \(197\), and
the second no factor \(53\). They already give a nonempty
different-owner packet.

Larger prescribed finite depths do not require prime values of a
polynomial. At \(k=1\), the derivative of the \(c\) factor is \(135\),
a unit modulo \(53\); the other arms are units. At \(k=2\), the
derivative of the \(A\) factor is \(315\), a unit modulo \(197\),
and again the other arms are units. For either simple root, the
expansion
\[
 F(r+q^j u)\equiv F(r)+q^j uF'(r)\pmod{q^{j+1}}
\]
has exactly one digit retaining the root. Lift to depth \(e\), then
choose a different next digit; the resulting residue has exact
depth \(e\). This realizes every prescribed positive integer
depth. CRT can additionally require the \(53\)-owner to be
\(1\pmod{197}\) and the \(197\)-owner to be \(2\pmod{53}\),
which excludes the other marked prime at each owner. Positive
representatives exist. For example depths four and five give
the marked cost \(\log53+2\log197>0\).
Theorem~\ref{cu19-density} supplies infinitely many actual pure
cubes retaining that fixed cost and the full new negative support.

The constant in \eqref{cu19-bb} depends on the fixed polynomial
\(G\), whose coefficients depend on the whole packet and its CRT
representative. The good root indices have density
\(d(\mathcal G)/H>0\) among positive integers, with no uniform
threshold for moving packets. The residual coefficient is one
and the exponent remains three, so the earlier convention
\(\lambda=1/n\) gives the fixed value \(1/3\).

These are fixed low-degree nonlinear selectors. The result does
not control the original \(n>324\) prime, \(B=n^4\) block, nor
allow the squarefree theorem to be invoked at an arbitrary changing
exponent without checking its irreducible factors and constants.
The elementary literal-cube identities have a separate scoped
Lean verification; CRT, squarefree density and its analytic input
remain ordinary proofs here. Other roots, general residuals,
exceptional sets, the independent-domain complement and full ABC
coverage remain open.
